\section*{RINGKASAN}
\addcontentsline{toc}{section}{RINGKASAN}

Perkembangan teknologi \textit{Artificial Intelligence} (AI) dan \textit{Machine Learning} (ML) menuntut penguatan kompetensi baru dalam pendidikan vokasional, khususnya bidang Rekayasa Perangkat Lunak. Siswa SMK tidak hanya perlu menguasai dasar pemrograman, tetapi juga membutuhkan pengalaman praktis membangun solusi berbasis data yang relevan dengan kebutuhan industri digital. Mitra kegiatan Pengabdian kepada Masyarakat ini adalah SMKN 2 Cimahi, Jurusan Rekayasa Perangkat Lunak, dengan permasalahan utama berupa terbatasnya akses pelatihan praktis \textit{Machine Learning} dan \textit{Computer Vision}, belum adanya pengalaman pembelajaran berbasis proyek kecerdasan buatan secara, serta belum optimalnya portofolio siswa untuk mendukung kesiapan kerja maupun studi lanjut.

Untuk menjawab permasalahan tersebut, kegiatan ini menyelenggarakan Workshop \textit{Machine Learning} untuk Prediksi Awal Kesehatan Tanaman melalui Analisis Citra Daun. Studi kasus dipilih karena bersifat visual, menggunakan \textit{dataset} publik gratis, tidak memerlukan perangkat keras khusus, dan relevan dengan kehidupan sehari-hari. Sistem diposisikan sebagai alat \textit{screening} awal, bukan diagnosis final, agar siswa memahami keterbatasan model dan pentingnya interpretasi hasil secara kritis. Pelaksanaan menggunakan pendekatan \textit{Project-Based Learning} yang mencakup penyusunan modul pelatihan, pengenalan konsep kecerdasan buatan, eksplorasi \textit{dataset} citra daun, pelatihan model klasifikasi menggunakan \textit{Google Colab}, evaluasi model, hingga pengembangan aplikasi prediksi sederhana berbasis \textit{Streamlit}. Evaluasi kegiatan dilakukan melalui \textit{pre-test} dan \textit{post-test}, penilaian \textit{mini project} kelompok, serta kuesioner kepuasan peserta, dengan dukungan mahasiswa ULBI sebagai fasilitator teknis.

Kegiatan ini diharapkan meningkatkan pemahaman dan keterampilan siswa dalam menerapkan \textit{Machine Learning} dan \textit{Computer Vision} secara praktis, menghasilkan bahan ajar yang dapat digunakan kembali oleh mitra, serta membentuk pengalaman belajar berbasis proyek yang memperkuat kompetensi digital siswa SMK.

\noindent\textbf{Kata Kunci} : \textit{Machine Learning}, \textit{Computer Vision}, \textit{Project-Based Learning}.
